
We have introduced relational graph convolutional networks (R-GCNs) and demonstrated their effectiveness in the context of two standard statistical relation modeling problems: link prediction and entity classification. For the entity classification problem, we have demonstrated that the R-GCN model can act as a competitive, end-to-end trainable graph-based encoder. For link prediction, the R-GCN model with DistMult factorization as the decoding component outperformed direct optimization of the factorization model, and achieved competitive results on standard link prediction benchmarks. Enriching the factorization model with an R-GCN encoder proved especially valuable for the challenging FB15k-237 dataset, yielding a 29.8\% improvement over the decoder-only baseline.

There are several ways in which our work could be extended. For example, the graph autoencoder model could be considered in combination with other factorization models, such as ComplEx~\cite{complex-complex_embeddings_for_simple_link_prediction}, which can be better suited for modeling asymmetric relations. It is also straightforward to integrate entity features in R-GCNs, which would be beneficial both for link prediction and entity classification problems. To address the scalability of our method, it would be worthwhile to explore subsampling techniques, such as in \citet{hamilton2017inductive}. Lastly, it would be promising to replace the current form of summation over neighboring nodes and relation types with a data-dependent attention mechanism. Beyond modeling knowledge bases, R-GCNs can be generalized to other applications where relation factorization models have been shown effective (e.g. relation extraction). 
